We now derive the time discretisations of \cref{eq:threespeciesmodel,eq:twospeciesmodel}. For both models, we use an IMEX scheme where we choose implicit discretisation for all linear terms and explicit discretisation for the non-linear reaction. For \cref{eq:twospeciesmodel}, this results in the iteration

\begin{equation}\label{eq:2s_weak_form_surface}
    (1 + \tau \gamma \delta_u) \int_{\Gamma_h} u_{n + 1} \psi \d S + \tau D_u \int_{\Gamma_h} \nabla u_{n + 1} \cdot \nabla \psi \d S = \int_{\Gamma_h} u_n \psi \d S + \tau \gamma \int_{\Gamma_h} \left[ P_u + \beta u_n^2 v_n \right] \psi \d S
\end{equation}

\noindent
for the surface PDE, where we now let $\Gamma_h$ denote the submesh\todo{define what a submesh is or use other notation in previous section} of $\partial \Omega_h$ that approximates $\Gamma$, and

\begin{equation}\label{eq:2s_weak_form_bulk}
    (1 + \tau \gamma \delta_v) \int_{\Omega_h} v_{n + 1} \phi + \tau D_v \int_{\Omega_h} \nabla v_{n + 1} \cdot \nabla \phi = \int_{\Omega_h} v_n \phi \d S - \tau \gamma \int_{\Gamma_h} \left[ P_u + \beta u_n^2 v_n \right] \phi \d S
\end{equation}

\noindent
for the bulk PDE, where $\psi \in \mathcal{V}_{\Gamma_h}$ and $\phi \in \mathcal{V}_{\Omega_h}$. With the same approach, \cref{eq:threespeciesmodel} admits the time discretisation

\begin{equation}\label{eq:3s_u}
    \begin{aligned}
        (1 + \tau \gamma \delta_u) \int_{\Gamma_h} u_{n + 1} \psi \d S + \tau D_u \int_{\Gamma_h} \nabla u_{n + 1} &\psi \d S - \tau \gamma (\alpha + 2 \koff) \int_{\Gamma_h} u_{b, n + 1} \psi \d S \\
        &= \int_{\Gamma_h} u_n \psi \d S + \tau \gamma \int_{\Gamma_h} \left[ P_u - 2 \kon u_n^2 v_n \right] \psi \d S
    \end{aligned}
\end{equation}

\noindent
for the PDE for $u$,

\begin{equation}\label{eq:3s_ub}
    \begin{aligned}
        (1 + \tau \kappa \gamma \delta_{u_b}) \int_{\Gamma_h} u_{b, n + 1} \chi \d S + \tau D_{u_b} \int_{\Gamma_h} \nabla u_{b, n + 1} &\chi \d S + \tau \kappa \gamma \koff \int_{\Gamma_h} u_{b, n + 1} \chi \d S \\
        &= \int_{\Gamma_h} u_{b,n} \chi \d S + \tau \kappa \gamma \kon \int_{\Gamma_h} u_n^2 v_n \chi \d S
    \end{aligned}
\end{equation}

\noindent
for the PDE for $u_b$, and finally for $v$,

\begin{equation}\label{eq:3s_v}
    \begin{aligned}
        (1 + \tau \gamma \delta_v) \int_{\Omega_h} v_{n + 1} \phi + \tau D_v \int_{\Omega_h} \nabla v_{n + 1} &\phi - \tau \gamma \koff \int_{\Gamma_h} u_{b, n + 1} \phi \d S \\
        &= \int_{\Omega_h} v_n \phi + \tau \gamma \int_{\Gamma_h} \left[ P_v - \kon u_n^2 v_n \right] \phi \d S,
    \end{aligned}
\end{equation}

\noindent
where $\psi, \chi \in \mathcal{V}_{\Gamma_h}$ and $\phi \in \mathcal{V}_{\Omega_h}$.

The initial conditions are set as random perturbations of the uniform steady states\todo{specify exact uniform steady states}. That is, we define $u_0, u_{b,0} \in \mathcal{V}_{\Gamma_h}$ and $v_0 \in \mathcal{V}_{\Omega_h}$ by

\begin{equation*}
    u_0 = \sum_{i = 1}^{l} (U + a_i) \psi_i, \qquad u_{b, 0} = \sum_{i = 1}^{l} (U_b + b_i) \psi_i \qquad \text{and} \qquad v_0 = \sum_{i = 1}^{m} (V + c_i) \phi_i,
\end{equation*}

\noindent
where $U$, $U_b$ and $V$ are the uniform steady states for $u$, $u_b$ and $v$, respectively, and $a_1, \dots, a_l$, $b_1, \dots, b_l$ and $c_1, \dots, c_m$ are uniformly sampled from $[-0.1, 0.1]$.

To calculate these iterations, \cref{eq:2s_weak_form_bulk,eq:2s_weak_form_surface,eq:3s_u,eq:3s_ub,eq:3s_v} are given as input to a Python scripts using FEniCSx, along with a mesh generated by Gmsh. FEniCSx derives the system of linear equations, which is solved with PETSc. All scripts that were written and used for this thesis can be found on the public GitHub page \url{https://github.com/bregtvdl03/Bachelor_Thesis_Turing_patterns}.